\section{Related Work}
\label{sec:related}

\paragraph{Serverless benchmark suites.} SeBS defined the template this work
follows: a benchmark specification, a cross-provider platform, and an
evaluation methodology, under the principles of relevance, usability,
scientific evaluation and extensibility~\cite{sebs2021}. SeBS-Flow extended it
to function workflows and, more importantly for us, demonstrated what
justifies a new suite: an execution model the existing tests cannot
characterise~\cite{sebsflow2025}. FunctionBench contributed micro and
application workloads for FaaS~\cite{functionbench2019}, and ServiTrace argues
for fine-grained distributed tracing and end-to-end application measurement
rather than aggregate latency~\cite{servitrace2025}. All of them measure
stateless invocation. None measures an interactive stateful session, a
suspend-resume round trip, or an environment reset.

\paragraph{Sandbox and microVM infrastructure.} Firecracker established the
microVM as the isolation boundary for multi-tenant
code execution~\cite{firecracker2020}; the gVisor case study quantified what
syscall interception costs~\cite{gvisor2019}; SOCK and Catalyzer attacked
container and VM startup latency~\cite{sock2018,catalyzer2020}. This body of
work measures the runtime. We measure the service built on top of one, and
treat the runtime as unobservable for hosted providers.

\paragraph{Reinforcement-learning environment execution.} EnvPool shows that
environment throughput, not only policy computation, bounds RL training
throughput~\cite{envpool2022}. It optimises a local execution engine. The
question we ask is one layer down: how the same environment workload behaves
when its environments are cloud sandboxes from different providers.

\paragraph{Agent benchmarks.} SWE-bench and its successors evaluate what a
model can accomplish, using a sandbox as a fixture~\cite{swebench2024}. We
evaluate the fixture.

\paragraph{Industry measurement.} The Kubernetes agent-sandbox project
publishes a cluster performance assessment covering rapid burst, high-volume
ramp, steady-state churn, throughput and warm-pool
burst~\cite{k8sagentsandbox}; our cluster load profiles follow its shapes
directly. Vendor and community benchmarks of hosted sandbox providers also
exist; we treat them as a source of load patterns rather than as comparable
measurements, since their configuration and client location are generally not
published.
